% GENERATED FILE. Edit canonical lessons, then run npm run generate:books.
% canonical-source-hash: fnv1a64:3cc3cd9bcc98d5dc
% canonical-lessons: TE-C272-pandiri, TE-C272-kance, TE-C272-poda, TE-C272-kandam, TE-C272-rekka, TE-R272-first-pass-words-from-chapters-267-269, TE-R272-second-pass-words-from-chapters-270-272

\chapter{Trellis, Fence, Bush, Stem, Wing}
\label{ch:te-trellis-fence-bush-stem-wing}
\hlchaptermodality{272}

\begin{chapteropening}
\textbf{By the end of this chapter:} \emph{I can say a trellis, a fence, a bush, a stem and a wing.}

Complete the last lesson of chapter 272: I can say a trellis, a fence, a bush, a stem and a wing.
\end{chapteropening}

\section[\texorpdfstring{\te{పందిరి} (pandiri)}{పందిరి (pandiri)}]{\te{పందిరి} (pandiri) — a trellis; a festive canopy}

\label{lesson:TE-C272-pandiri}



\begin{quote}
Before the new one: say the Telugu for foam, then the Telugu for a hole.
\end{quote}

\subsection*{You'll want to know: \te{పందిరి}}
\textbf{\te{పందిరి}} — \emph{pandiri} — \textquotedblleft{}a trellis; a festive canopy\textquotedblright{}.

\begin{grammarlens}[title={What you've built}]
One more word for this chapter.
\end{grammarlens}

\subsection*{Your turn}
\begin{itemize}
\raggedright
  \item \emph{Say it:} \emph{pandiri}
  \item \emph{Say it:} \emph{pandiri}, once more
  \item \emph{Say it:} say \emph{pandiri}
  \item \emph{Recall:} say the Telugu for foam, then the Telugu for a hole, then say \emph{pandiri} again
\end{itemize}

\begin{tcolorbox}[breakable,colback=teal!4,colframe=teal!35!black,title={Before you move on}]
What does \te{పందిరి} mean? (\textquotedblleft{}A trellis; a festive canopy\textquotedblright{}.) Say it once more.
\end{tcolorbox}
\section[\texorpdfstring{\te{కంచె} (kañce)}{కంచె (kañce)}]{\te{కంచె} (kañce) — a fence}

\label{lesson:TE-C272-kance}



\begin{quote}
Before the new one: say the Telugu for a hole, then the Telugu for a trellis.
\end{quote}

\subsection*{You'll want to know: \te{కంచె}}
\textbf{\te{కంచె}} — \emph{kañce} — \textquotedblleft{}a fence\textquotedblright{}.

\begin{grammarlens}[title={What you've built}]
One more word for this chapter.
\end{grammarlens}

\subsection*{Your turn}
\begin{itemize}
\raggedright
  \item \emph{Say it:} \emph{kañce}
  \item \emph{Say it:} \emph{kañce}, once more
  \item \emph{Say it:} say \emph{kañce}
  \item \emph{Recall:} say the Telugu for a hole, then the Telugu for a trellis, then say \emph{kañce} again
\end{itemize}

\begin{tcolorbox}[breakable,colback=teal!4,colframe=teal!35!black,title={Before you move on}]
What does \te{కంచె} mean? (\textquotedblleft{}A fence\textquotedblright{}.) Say it once more.
\end{tcolorbox}
\section[\texorpdfstring{\te{పొద} (poda)}{పొద (poda)}]{\te{పొద} (poda) — a bush}

\label{lesson:TE-C272-poda}



\begin{quote}
Before the new one: say the Telugu for a trellis, then the Telugu for a fence.
\end{quote}

\subsection*{You'll want to know: \te{పొద}}
\textbf{\te{పొద}} — \emph{poda} — \textquotedblleft{}a bush\textquotedblright{}.

\begin{grammarlens}[title={What you've built}]
One more word for this chapter.
\end{grammarlens}

\subsection*{Your turn}
\begin{itemize}
\raggedright
  \item \emph{Say it:} \emph{poda}
  \item \emph{Say it:} \emph{poda}, once more
  \item \emph{Say it:} say \emph{poda}
  \item \emph{Recall:} say the Telugu for a trellis, then the Telugu for a fence, then say \emph{poda} again
\end{itemize}

\begin{tcolorbox}[breakable,colback=teal!4,colframe=teal!35!black,title={Before you move on}]
What does \te{పొద} mean? (\textquotedblleft{}A bush\textquotedblright{}.) Say it once more.
\end{tcolorbox}
\section[\texorpdfstring{\te{కాండం} (kāṇḍaṁ)}{కాండం (kāṇḍaṁ)}]{\te{కాండం} (kāṇḍaṁ) — a stem}

\label{lesson:TE-C272-kandam}



\begin{quote}
Before the new one: say the Telugu for a fence, then the Telugu for a bush.
\end{quote}

\subsection*{You'll want to know: \te{కాండం}}
\textbf{\te{కాండం}} — \emph{kāṇḍaṁ} — \textquotedblleft{}a stem\textquotedblright{}.

\begin{grammarlens}[title={What you've built}]
One more word for this chapter.
\end{grammarlens}

\subsection*{Your turn}
\begin{itemize}
\raggedright
  \item \emph{Say it:} \emph{kāṇḍaṁ}
  \item \emph{Say it:} \emph{kāṇḍaṁ}, once more
  \item \emph{Say it:} say \emph{kāṇḍaṁ}
  \item \emph{Recall:} say the Telugu for a fence, then the Telugu for a bush, then say \emph{kāṇḍaṁ} again
\end{itemize}

\begin{tcolorbox}[breakable,colback=teal!4,colframe=teal!35!black,title={Before you move on}]
What does \te{కాండం} mean? (\textquotedblleft{}A stem\textquotedblright{}.) Say it once more.
\end{tcolorbox}
\section[\texorpdfstring{\te{రెక్క} (rekka)}{రెక్క (rekka)}]{\te{రెక్క} (rekka) — a wing}

\label{lesson:TE-C272-rekka}



\begin{quote}
Before the new one: say the Telugu for a bush, then the Telugu for a stem.
\end{quote}

\subsection*{You'll want to know: \te{రెక్క}}
\textbf{\te{రెక్క}} — \emph{rekka} — \textquotedblleft{}a wing\textquotedblright{}.

\begin{grammarlens}[title={What you've built}]
One more word for this chapter.
\end{grammarlens}

\subsection*{Your turn}
\begin{itemize}
\raggedright
  \item \emph{Say it:} \emph{rekka}
  \item \emph{Say it:} \emph{rekka}, once more
  \item \emph{Say it:} say \emph{rekka}
  \item \emph{Recall:} say the Telugu for a bush, then the Telugu for a stem, then say \emph{rekka} again
\end{itemize}

\begin{tcolorbox}[breakable,colback=teal!4,colframe=teal!35!black,title={Before you move on}]
What does \te{రెక్క} mean? (\textquotedblleft{}A wing\textquotedblright{}.) Say it once more.
\end{tcolorbox}
\section[(dialogue)]{Words from chapters 267-269 — a first pass}

\label{lesson:TE-R272-first-pass-words-from-chapters-267-269}



\begin{quote}
Say the words for a stem and a wing.
\end{quote}

\subsection*{Your turn}
Say these aloud:

\begin{itemize}
\raggedright
  \item a rainbow, a flood, a drought, moonlight, darkness
  \item a show, the radio, a programme, the internet, a website
  \item a tin, an eating plate, a sieve, wood, brass
\end{itemize}

\begin{tcolorbox}[breakable,colback=teal!4,colframe=teal!35!black,title={Before you move on}]
A rainbow? (\textbf{\te{ఇంద్రధనుస్సు}}, \emph{indradhanussu}.) A flood? (\textbf{\te{వరద}}, \emph{varada}.) A drought? (\textbf{\te{కరువు}}, \emph{karuvu}.) Moonlight? (\textbf{\te{వెన్నెల}}, \emph{vennela}.) Darkness? (\textbf{\te{చీకటి}}, \emph{cīkaṭi}.) A show? (\textbf{\te{ప్రదర్శన}}, \emph{pradarśana}.) The radio? (\textbf{\te{రేడియో}}, \emph{rēḍiyō}.) A programme? (\textbf{\te{కార్యక్రమం}}, \emph{kāryakramaṁ}.) The internet? (\textbf{\te{ఇంటర్నెట్}}, \emph{iṇṭarneṭ}.) A website? (\textbf{\te{వెబ్సైట్}}, \emph{vebsaiṭ}.) A tin? (\textbf{\te{డబ్బా}}, \emph{ḍabbā}.) An eating plate? (\textbf{\te{కంచం}}, \emph{kañcaṁ}.) A sieve? (\textbf{\te{జల్లెడ}}, \emph{jalleḍa}.) Wood? (\textbf{\te{చెక్క}}, \emph{cekka}.) Brass? (\textbf{\te{ఇత్తడి}}, \emph{ittaḍi}.)
\end{tcolorbox}
\section[(dialogue)]{Words from chapters 270-272 — a second pass}

\label{lesson:TE-R272-second-pass-words-from-chapters-270-272}



\begin{quote}
Say the words for a saucepan and a flowerpot.
\end{quote}

\subsection*{Your turn}
Say these aloud:

\begin{itemize}
\raggedright
  \item a saucepan, a flowerpot, a griddle, a mortar, a pestle
  \item lime, a roof tile, plastic, foam, a hole
  \item a trellis, a fence, a bush, a stem, a wing
\end{itemize}

\begin{tcolorbox}[breakable,colback=teal!4,colframe=teal!35!black,title={Before you move on}]
A saucepan? (\textbf{\te{తపేలా}}, \emph{tapēlā}.) A flowerpot? (\textbf{\te{కుండీ}}, \emph{kuṇḍī}.) A griddle? (\textbf{\te{పెనం}}, \emph{penaṁ}.) A mortar? (\textbf{\te{రోలు}}, \emph{rōlu}.) A pestle? (\textbf{\te{రోకలి}}, \emph{rōkali}.) Lime? (\textbf{\te{సున్నం}}, \emph{sunnaṁ}.) A roof tile? (\textbf{\te{పెంకు}}, \emph{peṅku}.) Plastic? (\textbf{\te{ప్లాస్టిక్కు}}, \emph{plāsṭikku}.) Foam? (\textbf{\te{నురుగు}}, \emph{nurugu}.) A hole? (\textbf{\te{చిల్లు}}, \emph{cillu}.) A trellis? (\textbf{\te{పందిరి}}, \emph{pandiri}.) A fence? (\textbf{\te{కంచె}}, \emph{kañce}.) A bush? (\textbf{\te{పొద}}, \emph{poda}.) A stem? (\textbf{\te{కాండం}}, \emph{kāṇḍaṁ}.) A wing? (\textbf{\te{రెక్క}}, \emph{rekka}.)
\end{tcolorbox}
